% GENERATED FILE. Edit canonical lessons, then run npm run generate:books.
% canonical-source-hash: fnv1a64:c7e048ca0644f0d3
% canonical-lessons: RU-C298-nayezd, RU-C298-postradavshiy, RU-C298-razbivatsya, RU-C298-utselet, RU-C298-perevernutsya

\chapter{On the road}
\label{ch:ru-on-the-road}
\hlchaptermodality{298}

\begin{chapteropening}
\textbf{By the end of this chapter:} \emph{I can describe a road accident: someone hit by a car, the injured, crashing, coming out unhurt and overturning.}

Complete the last lesson of chapter 298: I can describe a road accident: someone hit by a car, the injured, crashing, coming out unhurt and overturning.
\end{chapteropening}

\section[\texorpdfstring{\ru{наезд} (nayézd)}{наезд (nayézd)}]{\ru{наезд} (nayézd) — a hit (by a vehicle)}

\label{lesson:RU-C298-nayezd}



\begin{quote}
Before the new one: say the Russian for to be hurt, then the Russian for to wound.
\end{quote}

\subsection*{You'll want to know: \ru{наезд}}
\textbf{\ru{наезд}} (\emph{nayézd}) — \textquotedblleft{}a hit (by a vehicle)\textquotedblright{}.

\begin{quote}
\ru{Вчера} \ru{на} \ru{углу} \ru{был} \ru{наезд} \ru{на} \ru{мальчика}.
\end{quote}

\begin{quote}
\emph{Vcherá na uglú byl nayézd na málchika.}
\end{quote}

\begin{quote}
Yesterday a boy was hit by a car on the corner.
\end{quote}

\emph{\ru{Наезд} \ru{на}} + accusative.

\begin{cousinweb}
\emph{\ru{наезд}} is built on \emph{\ru{наехать}} \textquotedblleft{}to run into\textquotedblright{}, from \emph{\ru{на}-} and \emph{\ru{ехать}} \textquotedblleft{}to go by vehicle\textquotedblright{}, inherited from Proto-Slavic \emph{ěxati} \textquotedblleft{}to travel\textquotedblright{} (EtymDB).
\end{cousinweb}

\begin{grammarlens}[title={What you've built}]
One more word for telling, explaining and describing at B1.
\end{grammarlens}

\subsection*{Your turn}
\begin{itemize}
\raggedright
  \item \emph{Say it:} \emph{nayézd}
  \item \emph{Say it:} \emph{nayézd}, once more
  \item \emph{Say it:} a sentence of your own with \emph{nayézd}
  \item \emph{Recall:} say the Russian for to be hurt, then the Russian for to wound, then say \emph{nayézd} again
\end{itemize}

\begin{tcolorbox}[breakable,colback=teal!4,colframe=teal!35!black,title={Before you move on}]
What does \ru{наезд} mean? (\textquotedblleft{}A hit (by a vehicle)\textquotedblright{}.) Say it once more.
\end{tcolorbox}
\section[\texorpdfstring{\ru{пострадавший} (postradávshiy)}{пострадавший (postradávshiy)}]{\ru{пострадавший} (postradávshiy) — an injured person, a casualty}

\label{lesson:RU-C298-postradavshiy}



\begin{quote}
Before the new one: say the Russian for to wound, then the Russian for a hit.
\end{quote}

\subsection*{You'll want to know: \ru{пострадавший}}
\textbf{\ru{пострадавший}} (\emph{postradávshiy}) — \textquotedblleft{}an injured person, a casualty\textquotedblright{}.

\begin{quote}
\ru{Пострадавшего} \ru{отвезли} \ru{в} \ru{больницу}.
\end{quote}

\begin{quote}
\emph{Postradávshevo otvezlí v bolnítsu.}
\end{quote}

\begin{quote}
The injured man was taken to hospital.
\end{quote}

A participle that declines like an adjective; a woman is \emph{\ru{пострадавшая}}.

\begin{cousinweb}
\emph{\ru{пострадавший}} is the past participle of \emph{\ru{пострадать}} \textquotedblleft{}to be hurt\textquotedblright{}, from \emph{\ru{по}-} and \emph{\ru{страдать}} \textquotedblleft{}to suffer\textquotedblright{}; EtymDB lists no parent.
\end{cousinweb}

\begin{grammarlens}[title={What you've built}]
One more word for telling, explaining and describing at B1.
\end{grammarlens}

\subsection*{Your turn}
\begin{itemize}
\raggedright
  \item \emph{Say it:} \emph{postradávshiy}
  \item \emph{Say it:} \emph{postradávshiy}, once more
  \item \emph{Say it:} a sentence of your own with \emph{postradávshiy}
  \item \emph{Recall:} say the Russian for to wound, then the Russian for a hit, then say \emph{postradávshiy} again
\end{itemize}

\begin{tcolorbox}[breakable,colback=teal!4,colframe=teal!35!black,title={Before you move on}]
What does \ru{пострадавший} mean? (\textquotedblleft{}An injured person, a casualty\textquotedblright{}.) Say it once more.
\end{tcolorbox}
\section[\texorpdfstring{\ru{разбиваться} (razbivátsya)}{разбиваться (razbivátsya)}]{\ru{разбиваться} (razbivátsya) — to crash; to break, to smash}

\label{lesson:RU-C298-razbivatsya}



\begin{quote}
Before the new one: say the Russian for a hit, then the Russian for an injured person.
\end{quote}

\subsection*{You'll want to know: \ru{разбиваться}}
\textbf{\ru{разбиваться}} (\emph{razbivátsya}) — \textquotedblleft{}to crash; to break, to smash\textquotedblright{}.

\begin{quote}
\ru{Чашка} \ru{упала} \ru{и} \ru{разбилась}.
\end{quote}

\begin{quote}
\emph{Cháshka upála i razbílas.}
\end{quote}

\begin{quote}
The cup fell and broke.
\end{quote}

Imperfective; perfective \emph{\ru{разбиться}} (\emph{razbítsya}), \emph{\ru{разобьётся}} (\emph{razobyótsya}). Of a car or a plane: to crash.

\begin{cousinweb}
\emph{\ru{разбиваться}} is built on \emph{\ru{раз}-} \textquotedblleft{}apart\textquotedblright{} and \emph{\ru{бить}} \textquotedblleft{}to beat\textquotedblright{}, which is inherited from Proto-Slavic \emph{biti} \textquotedblleft{}to beat\textquotedblright{}, from PIE \emph{b\textsuperscript{h}eyH-} (EtymDB).
\end{cousinweb}

\begin{grammarlens}[title={What you've built}]
One more word for telling, explaining and describing at B1.
\end{grammarlens}

\subsection*{Your turn}
\begin{itemize}
\raggedright
  \item \emph{Say it:} \emph{razbivátsya}
  \item \emph{Say it:} \emph{razbivátsya}, once more
  \item \emph{Say it:} a sentence of your own with \emph{razbivátsya}
  \item \emph{Recall:} say the Russian for a hit, then the Russian for an injured person, then say \emph{razbivátsya} again
\end{itemize}

\begin{tcolorbox}[breakable,colback=teal!4,colframe=teal!35!black,title={Before you move on}]
What does \ru{разбиваться} mean? (\textquotedblleft{}To crash; to break, to smash\textquotedblright{}.) Say it once more.
\end{tcolorbox}
\section[\texorpdfstring{\ru{уцелеть} (utselét)}{уцелеть (utselét)}]{\ru{уцелеть} (utselét) — to come out unhurt, to survive intact}

\label{lesson:RU-C298-utselet}



\begin{quote}
Before the new one: say the Russian for an injured person, then the Russian for to crash.
\end{quote}

\subsection*{You'll want to know: \ru{уцелеть}}
\textbf{\ru{уцелеть}} (\emph{utselét}) — \textquotedblleft{}to come out unhurt, to survive intact\textquotedblright{}.

\begin{quote}
\ru{В} \ru{аварии} \ru{все} \ru{уцелели}.
\end{quote}

\begin{quote}
\emph{V avárii vse utseléli.}
\end{quote}

\begin{quote}
Everyone came out of the accident unhurt.
\end{quote}

Perfective. \emph{\ru{Дом} \ru{уцелел}} — the house was left standing.

\begin{cousinweb}
\emph{\ru{уцелеть}} is built on \emph{\ru{у}-} and \emph{\ru{целый}} \textquotedblleft{}whole, unhurt\textquotedblright{} (EtymDB glosses \emph{\ru{целый}} \textquotedblleft{}whole, healthy\textquotedblright{}): to stay whole.
\end{cousinweb}

\begin{grammarlens}[title={What you've built}]
One more word for telling, explaining and describing at B1.
\end{grammarlens}

\subsection*{Your turn}
\begin{itemize}
\raggedright
  \item \emph{Say it:} \emph{utselét}
  \item \emph{Say it:} \emph{utselét}, once more
  \item \emph{Say it:} a sentence of your own with \emph{utselét}
  \item \emph{Recall:} say the Russian for an injured person, then the Russian for to crash, then say \emph{utselét} again
\end{itemize}

\begin{tcolorbox}[breakable,colback=teal!4,colframe=teal!35!black,title={Before you move on}]
What does \ru{уцелеть} mean? (\textquotedblleft{}To come out unhurt, to survive intact\textquotedblright{}.) Say it once more.
\end{tcolorbox}
\section[\texorpdfstring{\ru{перевернуться} (perevernútsya)}{перевернуться (perevernútsya)}]{\ru{перевернуться} (perevernútsya) — to overturn, to roll over}

\label{lesson:RU-C298-perevernutsya}



\begin{quote}
Before the new one: say the Russian for to crash, then the Russian for to come out unhurt.
\end{quote}

\subsection*{You'll want to know: \ru{перевернуться}}
\textbf{\ru{перевернуться}} (\emph{perevernútsya}) — \textquotedblleft{}to overturn, to roll over\textquotedblright{}.

\begin{quote}
\ru{Машина} \ru{перевернулась} \ru{на} \ru{дороге}.
\end{quote}

\begin{quote}
\emph{Mashína perevernúlas na doróge.}
\end{quote}

\begin{quote}
The car overturned on the road.
\end{quote}

Perfective; imperfective \emph{\ru{переворачиваться}} (\emph{perevoráchivatsya}). Also of a person in bed: to turn over.

\begin{cousinweb}
\emph{\ru{перевернуться}} is built on \emph{\ru{пере}-} \textquotedblleft{}over\textquotedblright{} and \emph{-\ru{вернуть}}, the one-off form of \emph{\ru{вертеть}} \textquotedblleft{}to turn\textquotedblright{} (EtymDB), inherited from Proto-Slavic \emph{vĭrtěti}, PIE \emph{wert-} \textquotedblleft{}to turn\textquotedblright{}; modern \emph{\ru{вернуть}} on its own means \textquotedblleft{}to give back\textquotedblright{}.
\end{cousinweb}

\begin{grammarlens}[title={What you've built}]
One more word for telling, explaining and describing at B1.
\end{grammarlens}

\subsection*{Your turn}
\begin{itemize}
\raggedright
  \item \emph{Say it:} \emph{perevernútsya}
  \item \emph{Say it:} \emph{perevernútsya}, once more
  \item \emph{Say it:} a sentence of your own with \emph{perevernútsya}
  \item \emph{Recall:} say the Russian for to crash, then the Russian for to come out unhurt, then say \emph{perevernútsya} again
\end{itemize}

\begin{tcolorbox}[breakable,colback=teal!4,colframe=teal!35!black,title={Before you move on}]
What does \ru{перевернуться} mean? (\textquotedblleft{}To overturn, to roll over\textquotedblright{}.) Say it once more.
\end{tcolorbox}
